\section*{Chapter \ref{ch:b1:labour-supply-and-saving} --- Labour Supply and Saving}
% Numbers recomputed in code/bachelor-1/08-labour-supply-and-saving/exercises/solutions.py (tested).

\begin{solution}{exo:b1:labour-supply-and-saving:reservation}
The \omterm{def:b1:preferences-and-utility:mrs}{MRS} of leisure for consumption is $(1/\ell)/(1/c) = c/\ell$. At the point of
no work it is $6000/3000 = 2$: the \omterm{def:b1:labour-supply-and-saving:reservation-wage}{reservation wage} is 2 dalens an hour.
\end{solution}

\begin{solution}{exo:b1:labour-supply-and-saving:read}
France: $1504/2351 = 0.640$, a fall of $36.0\%$, while pay rose
$51.17/6.39 = 8.0$ times. United States: $1797/2022 = 0.889$, a fall of $11.1\%$,
while pay rose $45.86/13.48 = 3.4$ times.
\end{solution}

\begin{solution}{exo:b1:labour-supply-and-saving:tax}
Before: $c = 20(3000 - \ell) + 2000$, from $(\ell, c) = (3000, 2000)$ to $(0, 62000)$.
After: $c = 14(3000 - \ell) + 2000$, from $(3000, 2000)$ to $(0, 44000)$: the line
pivots about the \omterm{def:b1:consumer-choice:endowment}{endowment} point, as for a cut in the wage to 14.
\end{solution}

\begin{solution}{exo:b1:labour-supply-and-saving:income-sign}
Negative: the wage rise makes the worker richer at her current hours, and a
richer worker buys more leisure, so she works less. The income term is
$h\,\partial\ell/\partial m$: it is proportional to the hours she sells, since
the rise adds $h$ to her income for each dalen of wage.
\end{solution}

\begin{solution}{exo:b1:labour-supply-and-saving:log}
\omterm{def:b1:labour-supply-and-saving:full-income}{Full income} $wT + y$ is split equally: $\ell = (wT + y)/(2w)$ and
$h = T/2 - y/(2w)$. At $w = 20$, $y = 6000$: $h = 1500 - 150 = 1350$ hours.
\end{solution}

\begin{solution}{exo:b1:labour-supply-and-saving:implicit}
$60\% + 20\% = 80\%$. Of 1,000 extra dalens she keeps $1000 \times 0.2 = 200$.
\end{solution}

\begin{solution}{exo:b1:labour-supply-and-saving:withdrawal}
The net wage on the first hours falls from $60\%$ to $40\%$ of the gross wage,
so the gross \omterm{def:b1:labour-supply-and-saving:reservation-wage}{reservation wage} rises: \omterm{def:b1:labour-supply-and-saving:margins}{participation} falls. On the
withdrawal range the net wage is lower (substitution: fewer hours) and the
worker is poorer at given hours (\omterm{def:b1:income-and-substitution-effects:effects}{income effect}: more hours); the sign is
ambiguous. The benefit now runs out at a lower level of earnings (the same
benefit divided by $0.6$ instead of $0.4$). Workers between the new and the
old end of the range lose the rest of their benefit and now face no
withdrawal on their last dalens: both the higher net wage and the lower
income push their hours up. Workers above the old range are unaffected.
\end{solution}

\begin{solution}{exo:b1:labour-supply-and-saving:two-period}
$W = 30000 + 20000/1.05 = 49047.62$; $c_1 = W/1.9 = 25814.54$;
$c_2 = 0.9 \times 1.05 \times 25814.54 = 24394.74$; saving $30000 - 25814.54 = 4185.46$.
\end{solution}

\begin{solution}{exo:b1:labour-supply-and-saving:women}
Japan: $76.5 - 58.0 = 18.5$ points; United States: $68.9 - 66.5 = 2.4$ points. The
model says \omterm{def:b1:labour-supply-and-saving:margins}{participation} rises with the net wage and falls with non-labour
income (including a partner's income); to use it one would need the change
in women's net wages in each country, in taxes on second earners, in the
cost of childcare (which acts like a fixed cost of working) and in partners'
incomes. Norms about women's work are outside the model.
\end{solution}

\begin{solution}{exo:b1:labour-supply-and-saving:slutsky}
Multiply $\partial\ell/\partial w = \partial\ell^c/\partial w + h\,\partial\ell/\partial m$ by
$w/\ell$: $\varepsilon_{\ell w} = \varepsilon^c_{\ell w} + (wh/m)\,\eta_\ell$, with
$\eta_\ell = (m/\ell)\,\partial\ell/\partial m$. Since $h = T - \ell$,
$\varepsilon_{hw} = -(\ell/h)\varepsilon_{\ell w}$, and the same holds for the
compensated elasticities. So
$\varepsilon_{hw} = \varepsilon^c_{hw} - (\ell/h)(wh/m)\eta_\ell = \varepsilon^c_{hw} - (w\ell/m)\,\eta_\ell$.
\end{solution}

\begin{solution}{exo:b1:labour-supply-and-saving:borrower}
$s = y_1 - \bigl(y_1 + y_2/(1 + r)\bigr)/(1 + \beta)$, so
$\partial s/\partial r = y_2/\bigl((1 + \beta)(1 + r)^2\bigr) > 0$ whenever $y_2 > 0$, for
lenders and borrowers alike. With logarithmic utility the substitution and
\omterm{def:b1:income-and-substitution-effects:effects}{income effects} of $r$ on $c_1$ cancel, and only the fall in the present value of
$y_2$ remains. With preferences that substitute less between periods than
the logarithm, the \omterm{def:b1:income-and-substitution-effects:effects}{income effect} of a higher $r$ on a lender (she is richer)
can dominate, and she saves less.
\end{solution}

\begin{solution}{exo:b1:labour-supply-and-saving:refute}
The claim predicts that, across countries and over time, hours move
opposite to effective marginal tax rates on labour. It would be refuted by
countries whose tax rates rose with no fall in hours, or by a set of
European countries with similar tax rates but very different hours (or
the reverse), once differences in holidays set by law, unions and family
policy are accounted for. Within Europe, a group of countries with the
highest tax rates and also high employment, because taxes pay for services
that make work easier (childcare), would count against it.
\end{solution}

\begin{solution}{pb:b1:labour-supply-and-saving:second-job}
\textbf{1.} No tax; benefit $9000 - 0.48E$: $c = E + 9000 - 0.48E = 9000 + 0.52E$.

\textbf{2.} $c = 9000 + 0.52E - 0.2(E - 12000) = 11400 + 0.32E$. The benefit
$9000 - 0.48E$ is zero at $E = 9000/0.48 = 18750$.

\textbf{3.} $c = E - 0.2(E - 12000) = 2400 + 0.8E$.

\textbf{4.} $48\%$, $68\%$ and $20\%$.

\textbf{5.} $E = 12000$ at $12000/15 = 800$ hours, with $c = 9000 + 6240 = 15240$;
$E = 18750$ at 1,250 hours, with $c = 2400 + 15000 = 17400$.

\textbf{6.} Earnings rise from 12,000 to 18,000; $c$ rises from 15,240 to
$11400 + 0.32 \times 18000 = 17160$. She keeps $1920$ of $6000$, that is $32\%$: an
implicit marginal rate of $68\%$.

\textbf{7.} No. At 1,250 hours the slope rises from $0.32 \times 15 = 4.8$ to
$0.8 \times 15 = 12$ dalens an hour: the segment joining a point just below and a
point well above the kink lies above the \omterm{def:b1:consumer-choice:budget}{budget line}.

\textbf{8.} Net wage $0.52 \times 15 = 7.8$; \omterm{def:b1:labour-supply-and-saving:full-income}{full income} $9000 + 7.8 \times 3000 = 32400$;
$\ell = 0.6 \times 32400/7.8 = 2492.3$, $h = 507.7$, $c = 0.4 \times 32400 = 12960$.

\textbf{9.} $E = 15 \times 507.7 = 7615.5 < 12000$: on the first range. Utility
$0.4\ln 12960 + 0.6\ln 2492.3 = 8.4804$. On the third range the line's optimum
would be $\ell = 0.6 \times 38400/12 = 1920$, $h = 1080$, below the range, so her best
point there is the kink: $0.4\ln 17400 + 0.6\ln 1750 = 8.3861$, lower.

\textbf{10.} First range: $\ell = 0.5 \times 32400/7.8 = 2076.9$, $h = 923.1$, above 800:
off the range. Second: \omterm{def:b1:labour-supply-and-saving:full-income}{full income} $11400 + 4.8 \times 3000 = 25800$,
$\ell = 0.5 \times 25800/4.8 = 2687.5$, $h = 312.5$, below 800: off. Third: \omterm{def:b1:labour-supply-and-saving:full-income}{full income}
$2400 + 12 \times 3000 = 38400$, $\ell = 1600$, $h = 1400$, on the range.

\textbf{11.} With $a = 0.5$, utility is $\tfrac12\ln(c\ell)$. At the kink,
$15240 \times 2200 = 33528000$; at 1,400 hours, $c = 2400 + 0.8 \times 21000 = 19200$ and
$19200 \times 1600 = 30720000$. He works 800 hours, at the kink.

\textbf{12.} At $(c, \ell) = (9000, 3000)$ the \omterm{def:b1:preferences-and-utility:mrs}{MRS} is $\tfrac{a}{1 - a}\cdot\tfrac{c}{\ell}$: $1.5 \times 3 = 4.5$
for Ana and $3$ for Ben, in net dalens an hour. Divided by $0.52$: gross
\omterm{def:b1:labour-supply-and-saving:reservation-wage}{reservation wages} of $8.65$ and $5.77$.

\textbf{13.} A spike of workers with earnings of exactly 12,000 dalens, where
the implicit marginal rate jumps from $48\%$ to $68\%$.

\textbf{14.} Net wage $9.36$, \omterm{def:b1:labour-supply-and-saving:full-income}{full income} $9000 + 9.36 \times 3000 = 37080$,
$\ell = 0.6 \times 37080/9.36 = 2376.9$, $h = 623.1$ (earnings 11,216, still on the
range): her hours rise with the wage, because she has non-labour income.

\textbf{15.} Up to 12,000: $c = 9000 + 0.7E$, rate $30\%$. From 12,000 until the
benefit runs out at $9000/0.3 = 30000$: $c = 11400 + 0.5E$, rate $50\%$. From 30,000
to 40,000: $c = 2400 + 0.8E$, rate $20\%$.

\textbf{16.} Net wage $10.5$, \omterm{def:b1:labour-supply-and-saving:full-income}{full income} $9000 + 31500 = 40500$,
$\ell = 0.6 \times 40500/10.5 = 2314.3$, $h = 685.7$, $c = 16200$; earnings 10,285.7, on
the first range.

\textbf{17.} The net wage rose from 7.8 to 10.5. The \omterm{def:b1:income-and-substitution-effects:effects}{substitution effect} raises
hours; the \omterm{def:b1:income-and-substitution-effects:effects}{income effect}, since she is richer at any hours, lowers them.
Her hours rise by 178: the \omterm{def:b1:income-and-substitution-effects:effects}{substitution effect} wins.

\textbf{18.} Before: $9000 - 0.48 \times 7615.5 = 5344.6$. After:
$9000 - 0.30 \times 10285.7 = 5914.3$. Net cost $569.7$ dalens a year.

\textbf{19.} At 12,000 dalens of earnings his benefit goes from
$9000 - 5760 = 3240$ to $9000 - 3600 = 5400$: 2,160 dalens more, with no change in
hours.

\textbf{20.} Much of the cost goes to workers who do not change what they do:
those like Ben, and those earning between 18,750 and 30,000 who now receive a
benefit they did not get before. The extra hours of workers like Ana pay
back only a small part.

\textbf{21.} $c_1 + c_2/1.04 = 16200 + 12960/1.04$.

\textbf{22.} $W = 16200 + 12461.54 = 28661.54$; $c_1 = 28661.54/1.96 = 14623.23$;
$c_2 = 0.96 \times 1.04 \times 14623.23 = 14599.84$; saving $16200 - 14623.23 = 1576.77$.

\textbf{23.} $W = 16200 + 12960/1.08 = 28200$, $c_1 = 28200/1.96 = 14387.76$ and saving
1,812.24. Saving rises: with logarithmic utility $r$ moves $c_1$ only through the
present value of next year's income, which falls.

\textbf{24.} With $y_2 = 0$, $W = y_1$ and $c_1 = y_1/(1 + \beta)$, whatever $r$.

\textbf{25.} $W = 9000 + 30000/1.04 = 37846.15$, and he would like
$c_1 = 37846.15/1.96 = 19309.26$, borrowing 10,309.26. If he cannot borrow, he is at
the corner $c_1 = 9000$, $c_2 = 30000$.

\textbf{26.} On the second job she keeps 32 of every 100 dalens: the implicit
marginal rate on the second job is $68\%$.
\end{solution}
